\section{Instrument Design}
\label{sec:instrument-design}

HunStat2 is a single-channel, three-electrode potentiostat built directly
from the Analog Devices AD5940/AD5941 datasheet~\cite{devices2024ad5940},
standard three-electrode potentiostat design practice, and the Analog
Devices AD5940 evaluation-kit reference firmware~\cite{ad5940_eis_github}.
This section fixes the terms used throughout the rest of the paper: what
the instrument is built from (Section~\ref{sec:components}), how those parts
are wired into a measurement loop (Section~\ref{sec:block-diagram}), and the
nominal-component ``ideal'' values the dummy-cell measurements of
Section~\ref{sec:results} are compared against
(Section~\ref{sec:ideal-values}).

\subsection{Components}
\label{sec:components}
Table~\ref{tab:components} lists the instrument's principal components. The
AD5941 supplies the entire analog measurement chain -- excitation DACs,
potentiostat and transimpedance amplifiers, ADC, and DFT engine -- as
characterized, off-the-shelf silicon~\cite{devices2024ad5940}; the RP2040
MCU is responsible only for register sequencing, host communication, and
unit conversion, which is why Section~\ref{sec:intro} scopes the paper's
reliability question to firmware structure rather than analog design.

\begin{table}[h]
\centering
\caption{Principal components of the HunStat2 instrument.}
\label{tab:components}
\begin{threeparttable}
\begin{tabular}{@{}p{0.24\linewidth}p{0.68\linewidth}@{}}
\toprule
Component & Role / key specification \\
\midrule
Electrochemical AFE & Analog Devices AD5941: integrated LP/HS
  potentiostat loops, 12-bit LPDAC + 6-bit $V_{\mathrm{zero}}$ DAC, 12-bit
  HSDAC, LP/HS transimpedance amplifiers with a switchable feedback
  resistor ($\Rtia$: \SI{200}{\ohm}--\SI{160}{\kilo\ohm}, 8 steps), 16-bit
  $\Sigma\Delta$ ADC, and a hardware DFT accelerator~\cite{devices2024ad5940} \\
Microcontroller & Seeed Studio XIAO RP2040 (dual-core Cortex-M0+); runs
  the register-sequencing firmware and the USB-CDC host interface \\
MCU--AFE link & Bit-banged SPI (custom XIAO pin map; native SPI0 pinout
  not used), effective clock in the hundreds of kHz, well under the
  AD5941's rated \SI{12}{\mega\hertz}~\cite{devices2024ad5940} \\
Calibration resistor bank ($\Rcal$) & Jumper-selectable \SI{200}{\ohm},
  \SI{4.7}{\kilo\ohm}, \SI{10}{\kilo\ohm} on the carrier board; the
  reference the AFE measures against to calibrate its own transimpedance
  gain \\
Host interface & USB-CDC virtual serial, ASCII single-letter command
  protocol, \SI{1}{\mega baud} \\
Electrode/load interface & 3-terminal working/reference/counter (WE/RE/CE)
  connector; accepts either an electrochemical cell or the passive dummy
  cell of Section~\ref{sec:ideal-values} \\
\bottomrule
\end{tabular}
\end{threeparttable}
\end{table}

\subsection{Signal-chain / block diagram}
\label{sec:block-diagram}
Figure~\ref{fig:block-diagram} shows the instrument's signal chain, from
host command to applied excitation and back to reported current. The MCU
never touches the analog signal; every excitation, sensing, and gain
decision is made inside the AD5941, and the MCU's role is strictly
digital -- issue a register sequence, read back a code, convert it to
physical units. This division is what makes the firmware-architecture audit
of Section~\ref{sec:methods} meaningful in isolation: a bug in the firmware
layer cannot masquerade as, or be masked by, an analog defect, because the
two are architecturally separated at the SPI boundary in
Figure~\ref{fig:block-diagram}.

\begin{figure}[h]
\centering
\resizebox{\linewidth}{!}{%
\begin{tikzpicture}[
  block/.style={draw, rounded corners, minimum height=1.1cm,
    minimum width=2.3cm, align=center, font=\small},
  arr/.style={-{Latex[length=2mm]}, thick},
  edgelabel/.style={font=\scriptsize, align=center, text width=2.3cm}
]
\node[block] (pc) {Host PC\\(Python GUI / scripts)};
\node[block, right=2.7cm of pc] (mcu) {XIAO RP2040\\MCU};
\node[block, right=2.7cm of mcu] (afe) {AD5941\\AFE};
\node[block, right=2.9cm of afe] (cell) {3-electrode cell /\\dummy cell};

\draw[arr] (pc) to node[midway, above, edgelabel] {USB-CDC, \SI{1}{\mega baud}\\ASCII commands} (mcu);
\draw[arr] (mcu) to node[midway, above, edgelabel] {bit-banged SPI\\register read/write} (afe);
\draw[arr] ([yshift=2.2mm]afe.east) -- node[midway, above, edgelabel] {excitation\\(LPDAC/HSDAC $\to$ WE/CE)} ([yshift=2.2mm]cell.west);
\draw[arr] ([yshift=-2.2mm]cell.west) -- node[midway, below, edgelabel] {sense current\\(LPTIA/HSTIA $\to$ PGA $\to$ ADC)} ([yshift=-2.2mm]afe.east);
\end{tikzpicture}%
}
\caption{HunStat2 signal chain. The MCU sequences AD5941 registers over a
bit-banged SPI link and never handles the analog signal directly; the
AD5941 alone drives the excitation DACs, senses working-electrode current
through its transimpedance/PGA/ADC path, and reduces it to a raw code. The
host issues single-letter ASCII commands and receives method-tagged
data lines over USB-CDC.}
\label{fig:block-diagram}
\end{figure}

\subsection{Ideal dummy-cell reference values}
\label{sec:ideal-values}
The dummy-cell protocol of Section~\ref{sec:protocol} uses two passive
loads in place of an electrochemical cell: a Randles network
($R_s=\SI{560}{\ohm}$, $R_{ct}=\SI{10}{\kilo\ohm}$,
$C_{dl}=\SI{33}{\nano\farad}$) and a plain \SI{10}{\kilo\ohm} resistor. Both
have a current-voltage relationship that is fixed by their nominal
component values, independent of any potentiostat, which is what makes them
usable as an \emph{ideal} reference rather than merely a repeatable one. For
the DC methods reported here, the relevant zero-frequency resistance and
capacitive-loop terms are
\begin{align}
R_{\mathrm{DC}} &= R_s + R_{ct}, \label{eq:idc-rdc}\\
C_{\mathrm{eff}} &= \frac{R_{ct}^{2}\,C_{dl}}{(R_s+R_{ct})^{2}}, \label{eq:idc-ceff}
\end{align}
following the same Randles-network derivation used in the companion
metrology manuscript~\cite{clement2026linearity}. Under a linear potential
ramp at rate $\nu$, the ideal forward/reverse CV slope is $1/R_{\mathrm{DC}}$
and the ideal separation between the two legs is $\Delta I = 2\nu
C_{\mathrm{eff}}$. Table~\ref{tab:ideal} reports the resulting nominal
values for both loads at the $\nu=\SI{100}{\milli\volt\per\second}$ scan
rate used in Section~\ref{sec:protocol}.

\begin{table}[h]
\centering
\footnotesize
\setlength{\tabcolsep}{4pt}
\caption{Ideal (nominal-component) reference values for the two loads used
in the dummy-cell protocol (Section~\ref{sec:protocol}).}
\label{tab:ideal}
\begin{threeparttable}
\begin{tabular}{@{}p{4.0cm}p{1.9cm}p{1.9cm}@{}}
\toprule
Quantity & Randles cell & Resistor load \\
\midrule
$R_{\mathrm{DC}}$ (Eq.~\ref{eq:idc-rdc}) & \SI{10.56}{\kilo\ohm} & \SI{10.00}{\kilo\ohm} \\
$C_{\mathrm{eff}}$ (Eq.~\ref{eq:idc-ceff}) & \SI{29.6}{\nano\farad} & n/a (no capacitive branch) \\
Ideal CV slope $1/R_{\mathrm{DC}}$ & \SI{0.09470}{\micro\ampere\per\milli\volt} & \SI{0.10000}{\micro\ampere\per\milli\volt} \\
Ideal slope ratio (Randles:resistor) & \multicolumn{2}{p{4.1cm}}{0.9470} \\
Ideal CV loop separation $\Delta I$ at \SI{100}{\milli\volt\per\second} & \SI{5.9}{\nano\ampere} (sub-LSB) & 0 (purely resistive) \\
Ideal CA/SWV/DPV response & \multicolumn{2}{p{4.1cm}}{Flat plateau at $\Delta V/R_{\mathrm{DC}}$; no Faradaic peak (Section~\ref{sec:outlook-real-solution})} \\
\bottomrule
\end{tabular}
\begin{tablenotes}\footnotesize
\item Both loads are passive $R$/$RC$ networks with no charge-transfer
  reaction, so no method is expected to show a Faradaic peak; CA is
  expected to settle to a flat plateau (no Cottrell decay) and SWV/DPV are
  expected to report a differential current at or below the instrument's
  noise floor. Section~\ref{sec:results} compares the instrument's measured
  output against these values; the CV slope predictions above match the
  measured slopes of Section~\ref{sec:results} to within
  \SI{0.9}{\percent}, and the predicted sub-LSB loop separation is
  consistent with the visually overlapping forward/reverse legs reported
  there.
\end{tablenotes}
\end{threeparttable}
\end{table}
